% ==========================================================
\section{问题 1：菌体生长动力学模型}

\subsection{模型构建}
记 $L(t)=\lg N(t)$，其中 $N$ 为每毫升活菌数。由于菌体在迟滞期、对数期、稳定期的增长方式不同，本文在同一对数尺度上并列考察五类模型（表 \ref{tab:q1models}），它们的区别在于“限制方式”与“对称性”。

\begin{table}[H]
\centering
\caption{五类候选生长模型}
\label{tab:q1models}
\small
\begin{tabular}{llc}
\toprule
编号 & 模型 & 参数个数 \\
\midrule
M1 & 指数（Malthus）：$L=L_0+\mu t$ & 2 \\
M2 & 含迟滞期指数：$L=L_0+\mu\max(0,t-\lambda)$ & 3 \\
M3 & Logistic：$L=L_0+\dfrac{A}{1+e^{-k(t-t_m)}}$ & 4 \\
M4 & Gompertz：$L=L_0+A\exp\{-\exp[1+\frac{\mu_m e}{A}(\lambda-t)]\}$ & 4 \\
M5 & Richards：$L=L_0+A(1+\nu e^{-k(t-t_m)})^{-1/\nu}$ & 5 \\
\bottomrule
\end{tabular}
\end{table}

五类模型有共同的生物学含义：
$L_0$ 是迟滞期的水平，$A$ 是对数增量的上限，$t_m$（或 $\lambda$）刻画拐点位置，$k$（或 $\mu_m$）刻画增长快慢。
M3 假设“增长速度与剩余增长空间成正比”，曲线关于拐点对称；M4 假设“增长速度按指数衰减”，曲线左偏、有明显迟滞期；M5 用形状参数 $\nu$ 把二者统一（$\nu\to0$ 趋于 M4，$\nu=1$ 即 M3），因此 M5 可用来检验“对称性假设是否成立”。

\subsection{参数估计方法}
对每个模型求解加权非线性最小二乘问题
$$\hat{\boldsymbol\theta}=\arg\min_{\boldsymbol\theta}\sum_{i=1}^{n}\big[L_i-L(t_i;\boldsymbol\theta)\big]^2 ,$$
采用信赖域反射算法（Trust-Region-Reflective）迭代，并取 8 组不同的初值做多起点计算以避免落入局部极小；收敛判据取梯度、步长与目标函数变化的相对容差均为 $10^{-13}$。参数的渐近标准误由
$$\widehat{\mathrm{Cov}}(\hat{\boldsymbol\theta})=\hat\sigma^2\big(J^{\top}J\big)^{-1},\qquad \hat\sigma^2=\frac{\mathrm{RSS}}{n-p}$$
给出，其中 $J$ 为残差对参数的雅可比矩阵。

\subsection{参数估计结果}
五个模型的拟合结果见表 \ref{tab:q1fit}，Logistic 模型的参数估计与置信区间见表 \ref{tab:q1par}。

\begin{table}[H]
\centering
\caption{五类模型的拟合优度比较（$n=8$）}
\label{tab:q1fit}
\small
\begin{tabular}{lcccccc}
\toprule
模型 & $R^2$ & RMSE & MAPE/\% & AIC & BIC & 留一 RMSE \\
\midrule
M1 指数            & 0.9214 & 0.2295 & 2.11 & $-19.55$ & $-19.39$ & 0.3156 \\
M2 含迟滞期指数    & 0.9214 & 0.2295 & 2.11 & $-17.55$ & $-17.31$ & 0.3305 \\
M4 Gompertz        & 0.9983 & 0.0342 & 0.35 & $-46.02$ & $-45.70$ & 0.1014 \\
M5 Richards        & 0.99986 & 0.0099 & 0.10 & $-63.92$ & $-63.52$ & 0.0409 \\
\textbf{M3 Logistic} & \textbf{0.99982} & \textbf{0.0110} & \textbf{0.096} & $\mathbf{-64.12}$ & $\mathbf{-63.80}$ & \textbf{0.0273} \\
\bottomrule
\end{tabular}
\end{table}

\begin{table}[H]
\centering
\caption{Logistic 模型参数估计（标准误与 95\% 自助置信区间，$B=600$）}
\label{tab:q1par}
\small
\begin{tabular}{lcccc}
\toprule
参数 & 估计值 & 标准误 & 95\% 置信区间 & 含义 \\
\midrule
$L_0$ & 7.582 & 0.018 & $[7.556,\ 7.605]$ & 下渐近线（迟滞期水平） \\
$A$   & 2.106 & 0.024 & $[2.074,\ 2.135]$ & 对数增量 \\
$k$   & 0.683 $\mathrm{h^{-1}}$ & 0.021 & $[0.655,\ 0.713]$ & 速率常数 \\
$t_m$ & 5.511 h & 0.049 & $[5.440,\ 5.571]$ & 拐点时刻 \\
\bottomrule
\end{tabular}
\end{table}

于是菌体生长模型为
\begin{equation}
L(t)=7.582+\frac{2.106}{1+e^{-0.683(t-5.511)}}\qquad \big[\lg(\mathrm{CFU/mL})\big],
\label{eq:logistic}
\end{equation}
或写成等价的微分形式（更便于后续耦合）
\begin{equation}
\frac{\dd L}{\dd t}=\frac{k}{A}(L-L_0)(L_\infty-L),\qquad L_\infty=L_0+A=9.688 .
\label{eq:logistic-ode}
\end{equation}

\subsection{模型检验}

\paragraph{(1) 拟合精度。}决定系数 $R^2=0.99982$，RMSE 为 $0.0110\ \lg$ 单位，而原始数据只记录到 $0.01$ 一位小数，说明残差已降到四舍五入误差的量级；平均相对误差 $0.096\%$。逐点对比见表 \ref{tab:q1point}。

\begin{table}[H]
\centering
\caption{Logistic 模型逐点拟合结果}
\label{tab:q1point}
\small
\begin{tabular}{cccccc}
\toprule
$t$/h & 实测 $L$ & 拟合 $\hat L$ & 残差 & 相对误差/\% \\
\midrule
0  & 7.63 & 7.629 & $-0.001$ & $-0.01$ \\
2  & 7.75 & 7.757 & $+0.007$ & $+0.09$ \\
4  & 8.15 & 8.135 & $-0.015$ & $-0.19$ \\
6  & 8.79 & 8.809 & $+0.019$ & $+0.21$ \\
8  & 9.38 & 9.362 & $-0.018$ & $-0.19$ \\
10 & 9.59 & 9.594 & $+0.004$ & $+0.04$ \\
12 & 9.66 & 9.663 & $+0.003$ & $+0.03$ \\
14 & 9.68 & 9.681 & $+0.001$ & $+0.01$ \\
\bottomrule
\end{tabular}
\end{table}

\paragraph{(2) 残差结构。}残差在 $[-0.019,0.019]$ 内正负交替（图 \ref{fig:q1}(b)），无“先正后负”的系统性弯曲，说明模型没有遗漏趋势项；残差的绝对值随时间不增大，说明误差方差近似齐性，满足最小二乘的基本前提。

\paragraph{(3) 泛化能力。}留一交叉验证（每次去掉一个点重新拟合并预测该点）的 RMSE 为 $0.0273$，仅比拟合 RMSE 大 $0.016$，说明 4 个参数对 8 个点而言并不过拟合；作为对照，Richards 模型多 1 个参数，留一 RMSE 反而升至 $0.0409$，这是典型的“过拟合导致外推变差”。

\paragraph{(4) 模型选择。}按 AIC 准则，M3 的 AIC 为 $-64.12$，比 M5 低 $0.20$、比 M4 低 $18.1$、比 M1 低 $44.6$。M5 的形状参数估计为 $\hat\nu=1.28$，其 95\% 置信区间包含 $\nu=1$，即“Richards 模型退化为 Logistic 模型”，这从另一个角度证明对称 S 形是恰当的函数形式。综合精度、简约性与可解释性，本文选用 M3。

\paragraph{(5) 无模型对照。}不依赖任何函数形式，用中心差分直接估计斜率，最大值为 $0.308\ \lg\cdot\mathrm{h^{-1}}$（出现在 $t=5$ h 附近）；用三次平滑样条求导得 $0.274$。二者均小于模型值 $0.360$，原因是对 8 个离散点做差分/样条会低估峰值（差分是区间平均、样条带平滑惩罚），因此模型值 $0.360$ 更接近真实的瞬时最大速率，这与图 \ref{fig:q1}(c) 中“模型曲线包住差分点”的现象一致。

\begin{figure}[H]
\centering
\includegraphics[width=0.98\textwidth]{fig2_q1.png}
\caption{问题 1 的拟合结果。(a) 三种模型与实测值；(b) Logistic 模型残差；(c) 比生长速率曲线与差分估计对照}
\label{fig:q1}
\end{figure}

\subsection{结果与解读}

\paragraph{(1) 最大生长速率。}对式 \eqref{eq:logistic} 求导并令其最大，
$$\left.\frac{\dd L}{\dd t}\right|_{max}=\frac{kA}{4}=\frac{0.683\times2.106}{4}=0.360\ \ \lg(\mathrm{CFU/mL})\cdot\mathrm{h^{-1}},$$
出现在 $t=t_m=5.51$ h。自助法给出的 95\% 置信区间为 $[0.349,\ 0.371]$。
换算成自然对数的比生长速率：$\mu=\mu_m\ln 10=0.828\ \mathrm{h^{-1}}$，对应倍增时间 $t_d=\ln 2/\mu=0.84$ h。这一量级与乳酸菌在 37 ℃ 下的文献值相符，说明估计合理。

\paragraph{(2) 最终菌体量。}$t\to\infty$ 时 $L_\infty=L_0+A=9.688$，即 $N_\infty=10^{9.688}=4.87\times10^{9}$ CFU/mL；14 h 实测 $9.68$，对应 $4.79\times10^{9}$ CFU/mL，二者只差 $1.7\%$。相对接种时的 $7.63$（$4.27\times10^{7}$ CFU/mL），菌体量增长了约 $114$ 倍。

\paragraph{(3) 分期。}由 $\dd L/\dd t$ 曲线可定量划分三个阶段：$0\sim3$ h 速率小于峰值的 $20\%$，为迟滞期；$3\sim9$ h 速率大于峰值的 $50\%$，为对数期；$9$ h 之后速率降到峰值的 $10\%$ 以下，进入稳定期。这与图 \ref{fig:data}(a) 的目视判断一致，但给出了明确的边界，为问题 2 中“为什么后期产酸仍能继续”的解释提供了依据。

% ==========================================================
\section{问题 2：生长—产酸—耗糖耦合模型}

\subsection{联动关系的识别}
在建模之前，先用数据回答“三者到底怎样联动”。

\paragraph{(1) 产酸与耗糖。}以累计耗糖 $S_c=S_0-S$ 为横轴、乳酸 $P$ 为纵轴（图 \ref{fig:q2}(a)），8 个点几乎落在一条过原点的直线上，相关系数 $0.9996$。过原点最小二乘得
\begin{equation}
P=Y_{P/S}\,S_c,\qquad \hat Y_{P/S}=\frac{\sum P_iS_{c,i}}{\sum S_{c,i}^2}=0.9893\ \mathrm{g/g},
\label{eq:yield}
\end{equation}
$R^2=0.9988$、$\mathrm{RMSE}=0.133$ g/L。若允许截距，则 $P=1.0035S_c-0.1102$，截距项不显著，进一步支持“过原点”的化学计量结构。

这一结果的物理意义很强：同型乳酸发酵的反应式为
$$\mathrm{C_6H_{12}O_6\longrightarrow 2\,C_3H_6O_3},$$
葡萄糖与乳酸的摩尔质量均为 $180\ \mathrm{g/mol}$，故理论得率恰为 $1.0$ g/g。实测 $0.989$ 与理论值只差 $1.1\%$，说明本批发酵以同型乳酸发酵为主，几乎没有副产物（乙醇、乙酸）分流碳源。

\paragraph{(2) 产酸与生长。}把 $P$ 对“菌体对数增量” $\Delta L=L-L_0$ 做图可见明显的上凸：前 6 h 每增加 1 个数量级菌体只产 $2.1$ g/L 酸，后 6 h 则达 $4.5$ g/L。这说明产酸不能只用“与生长成正比”解释，必须有“与菌体量成正比”的维持项。这正是 Luedeking--Piret 方程的结构。

\subsection{模型构建}
综合以上两点，建立三方程耦合模型。以相对菌体量 $X=10^{L-L_0}$ 表示菌体多少（接种时刻 $X=1$），则
\begin{align}
\frac{\dd P}{\dd t}&=\Big(\alpha\frac{\dd L}{\dd t}+\beta X\Big)\Big(1-\frac{P}{P_m}\Big), \label{eq:lp}\\
\frac{\dd S}{\dd t}&=-\frac{1}{Y_{P/S}}\frac{\dd P}{\dd t}-m_s, \label{eq:sugar}\\
L(t)&=7.582+\frac{2.106}{1+e^{-0.683(t-5.511)}} \quad\text{（问题 1 的结果）} \label{eq:linput}
\end{align}
其中：
\begin{itemize}[leftmargin=2em,itemsep=0.08em]
\item $\alpha\,\dd L/\dd t$ 为\wj{生长偶联项}：菌体每增长一个数量级所伴随的产酸，反映“长菌就要产酸”；
\item $\beta X$ 为\wj{非生长偶联项}：单位菌体在单位时间内的基础产酸，反映“不长菌也能产酸”（维持代谢与静息细胞的产酸）；
\item $(1-P/P_m)$ 为\wj{产物抑制项}：乳酸累积使胞内 $H^+$ 与未解离乳酸升高，抑制乳酸脱氢酶活性，$P_m$ 为效价上限；
\item $1/Y_{P/S}$ 与 $m_s$ 把耗糖分解为“产酸耗糖”与“维持耗糖”。
\end{itemize}

\noindent 模型共 5 个待估参数 $\boldsymbol\theta=(\alpha,\beta,P_m,Y_{P/S},m_s)$，状态变量为 $(P,S)$，$L(t)$ 作为外部输入。这样处理有两个好处：一是与问题 1 严格衔接，避免了“为拟合 $L$ 而牺牲 $P$”的折中；二是把误差来源分离——$L$ 的误差已由问题 1 控制，问题 2 只需解释 $P$ 与 $S$。

\subsection{参数估计}
给定 $\boldsymbol\theta$，用 LSODA 变步长法积分式 \eqref{eq:lp}--\eqref{eq:sugar}（相对容差 $10^{-10}$，最大步长 $0.05$ h），得到 $\hat P(t_i),\hat S(t_i)$，再极小化加权残差
$$\min_{\boldsymbol\theta}\sum_{i=1}^{8}\left[\left(\frac{\hat P_i-P_i}{\sigma_P}\right)^2+\left(\frac{\hat S_i-S_i}{\sigma_S}\right)^2\right],\quad \sigma_P=\sigma_S=0.15\ \mathrm{g/L}.$$
采用信赖域反射算法，8 组初值多起点计算，参数估计与不确定度见表 \ref{tab:q2par}。标准误由雅可比矩阵给出，置信区间由 400 次残差自助法给出。

\begin{table}[H]
\centering
\caption{问题 2 模型参数估计}
\label{tab:q2par}
\small
\begin{tabular}{lcccc}
\toprule
参数 & 估计值 & 标准误 & 95\% 置信区间 & 物理含义 \\
\midrule
$\alpha$ & 2.269 & 0.408 & $[1.634,\ 2.750]$ & 生长偶联产酸系数 (g/L per lg) \\
$\beta$  & 0.0466 $\mathrm{h^{-1}}$ & 0.0087 & $[0.0358,\ 0.0611]$ & 非生长偶联产酸系数 \\
$P_m$    & 9.912 g/L & 0.340 & $[9.424,\ 10.482]$ & 产物抑制阈值 \\
$Y_{P/S}$& 1.056 & 0.187 & $[0.952,\ 1.355]$ & 乳酸对糖的得率 (g/g) \\
$m_s$    & 0.0459 g/(L$\cdot$h) & 0.121 & $[0,\ 0.200]$ & 维持耗糖速率 \\
\bottomrule
\end{tabular}
\end{table}

\subsection{拟合效果}

\paragraph{(1) 整体精度。}乳酸：$R^2=0.9955$、$\mathrm{RMSE}=0.257$ g/L；残糖：$R^2=0.9946$、$\mathrm{RMSE}=0.280$ g/L（图 \ref{fig:q2}(b)）。考虑到数据本身只有两位小数且取样间隔 2 h，这一精度已接近数据分辨极限。

\begin{table}[H]
\centering
\caption{问题 2 模型逐点拟合结果}
\label{tab:q2point}
\small
\begin{tabular}{ccccccccc}
\toprule
$t$/h & $P$实测 & $\hat P$ & 误差 & 相对\% & $S$实测 & $\hat S$ & 误差 & $P/S_c$ \\
\midrule
0  & 0.00 & 0.00 & 0.00 & --  & 35.00 & 35.00 & 0.00 & -- \\
2  & 0.36 & 0.40 & $+0.04$ & $+10.8$ & 34.60 & 34.53 & $-0.07$ & 0.900 \\
4  & 1.08 & 1.37 & $+0.29$ & $+26.7$ & 33.80 & 33.52 & $-0.28$ & 0.900 \\
6  & 2.88 & 3.14 & $+0.26$ & $+9.1$  & 31.80 & 31.75 & $-0.05$ & 0.900 \\
8  & 6.12 & 5.68 & $-0.44$ & $-7.3$  & 28.70 & 29.26 & $+0.56$ & 0.971 \\
10 & 7.83 & 8.08 & $+0.25$ & $+3.2$  & 27.10 & 26.88 & $-0.22$ & 0.991 \\
12 & 9.09 & 9.29 & $+0.20$ & $+2.2$  & 26.00 & 25.65 & $-0.35$ & 1.010 \\
14 & 9.99 & 9.72 & $-0.27$ & $-2.7$  & 24.90 & 25.15 & $+0.25$ & 0.989 \\
\bottomrule
\end{tabular}
\end{table}

\paragraph{(2) 残差的位置。}最大残差出现在 $t=8$ h（$-0.44$ g/L），此处正是菌体从对数期转入稳定期的过渡带；$4\sim6$ h 的相对误差偏大是因为分母本身很小（1 g/L 量级），绝对误差仅 $0.3$ g/L。残差符号在时间上交替出现，没有“持续偏高/偏低”的系统偏差。

\paragraph{(3) 与简化模型的比较。}若不考虑产物抑制（令 $P_m\to\infty$），模型退化为 4 参数，此时乳酸 $R^2$ 降至 $0.9883$，且预测的终点乳酸为 $10.23$ g/L（高于实测 $9.99$），AIC 由 $28.69$ 升至 $37.26$。若进一步去掉生长偶联项（只留 $\beta X$），$R^2$ 降至 $0.9819$，且早期产酸被严重低估。这说明“两个偶联项 + 产物抑制”三者缺一不可。

\paragraph{(4) 外推检验。}只用前 6 个点（$0\sim10$ h）重新估计参数，再外推 12 h 与 14 h，得到 $\hat P(12)=7.97$、$\hat P(14)=7.98$，明显低于实测的 $9.09$ 与 $9.99$。原因很清楚：$0\sim10$ h 内乳酸最高只到 $7.83$ g/L，尚未“触碰”抑制阈值 $P_m$，参数 $P_m$ 在此时不可辨识（估成 $7.98$），模型于是错误地认为平台在 8 g/L 就出现。\wj{这提醒我们：含阈值/饱和参数的模型，必须用跨越阈值的完整数据来定参，不能凭前段数据外推后段。}本文最终使用全部 8 个点，$P_m$ 的置信区间 $[9.42,10.48]$ 宽度仅 $1.1$ g/L，是可辨识的。

\paragraph{(5) 敏感性。}把每个参数单独放大 $10\%$，终点乳酸的变化率为：$P_m$ 使 $+9.06\%$、$\beta$ 使 $+0.58\%$、$\alpha$ 使 $+0.09\%$，$Y_{P/S}$ 与 $m_s$ 几乎无影响（它们只改变耗糖）。可见\wj{终点产量几乎完全由抑制阈值决定}，这是后文所有工艺建议的物理基础。

\begin{figure}[H]
\centering
\includegraphics[width=0.98\textwidth]{fig3_q2.png}
\caption{问题 2 的结果。(a) 产酸—耗糖化学计量关系；(b) 耦合模型对乳酸与残糖的拟合；(c) 分段产酸速率对比；(d) 菌体量的对数增长}
\label{fig:q2}
\end{figure}

\subsection{机理解读}

\paragraph{(1) 产酸来源分解。}把式 \eqref{eq:lp} 的两项分别沿轨迹积分，得生长偶联项贡献 $3.14$ g/L（$31.8\%$）、非生长偶联项贡献 $6.75$ g/L（$68.2\%$）。即\wj{本批发酵的乳酸约有三分之二来自“不长菌的维持代谢”}，这解释了为什么菌体进入稳定期后产酸仍在继续，也说明单纯追求“长得快”并不能等比例提高产量。

\paragraph{(2) 为什么只用了 10 g/L 糖。}模型给出的抑制阈值 $P_m=9.91$ g/L，而 14 h 时乳酸已达 $9.99$ g/L，$(1-P/P_m)\approx0$，产酸速率趋于零；同时残糖仍有 $24.9$ g/L。若继续发酵到 28 h，模型预测乳酸仅升至 $9.91$ g/L、残糖降至 $24.3$ g/L，即“再怎么等也没有意义”。这从机理上确认：\wj{本批发酵是被自己的产物“憋停”的，不是缺糖}。该结论直接引出问题 3 与问题 4 的分析方向。

\paragraph{(3) 得率的自洽性。}$Y_{P/S}=1.056$，其 95\% 置信区间 $[0.952,1.355]$ 包含理论值 $1.0$；而独立由式 \eqref{eq:yield} 得到的 $0.989$ 与之相差 $6\%$，两条路径互相印证。取二者平均，可认为\wj{每消耗 1 g 葡萄糖约生成 0.99--1.06 g 乳酸}。
